\section{Introduction}

Large language model (LLM) based agents have recently demonstrated strong performance across a wide range of complex tasks, including software engineering, scientific discovery, robotics planning, and data analysis~\cite{plaat2025agentic, jimenez2023swe, boiko2023emergent, singh2024malmm}. 
By combining instruction following, multi-step reasoning, tool use, and access to external knowledge, agent frameworks provide a flexible paradigm for solving problems that extend beyond single-turn prediction or classification. 
Motivated by these successes, a growing body of work has begun to explore whether agentic approaches can also benefit time-series applications~\cite{zhao2025timeseriesscientist, zhang2024large, yeh2025empowering}.

At the same time, the time-series community has developed a rich ecosystem of specialized forecasting models and foundation models that achieve state-of-the-art numerical accuracy across many domains~\cite{cao2025timedit, ansari2024chronos, yang2025foundation, cao2023tempo}. 
These advances raise a fundamental question: \textit{if forecasting performance is already strong, what additional capabilities, if any, do agents bring to time-series analysis and forecasting?} 
Do agent-based systems merely reproduce numerical predictions, or do they enable forecasting behavior that adapts to contextual information and changing conditions~\cite{chang2025survey, chow2024towards, liu2025ts}?

Modern forecasting models effectively capture temporal patterns such as trends, seasonality, and local dependencies, which are sufficient for accurate numerical prediction~\cite{song2024deep, liu2024deep}. 
However, this use of temporal structure is largely implicit and optimized for minimizing prediction error, rather than for producing representations that can be explicitly interpreted, manipulated, or combined with contextual information. 
As a result, strong numerical performance alone does not reveal which temporal competencies a model possesses, nor whether temporal signals can be reused under different informational settings. 
This motivates evaluation frameworks that isolate distinct aspects of temporal competence, rather than aggregating them into a single end-to-end forecasting task. 
As we later show, models with strong forecasting accuracy often fail to adapt their predictions under contextual or event-driven conditions.

A key distinction in real-world temporal decision-making lies in the role of \emph{context}. 
Contextual information fundamentally shapes how temporal patterns should be interpreted, rather than merely refining numerical estimates. 
In many applications, events, domain semantics, or external conditions determine how historical trends inform future predictions~\cite{williams2024context, wang2024news}. 
For example, in retail demand planning, knowledge of an upcoming heatwave may outweigh historical sales trends when forecasting demand. 
Evaluating such behavior requires tasks that selectively introduce context and events, allowing us to assess whether forecasting mechanisms generalize across different informational regimes.

Despite increasing interest in agent-based time-series analysis, existing benchmarks are not designed to isolate and measure these capabilities. 
Forecasting benchmarks primarily assess numerical prediction error, without evaluating whether models correctly interpret historical structure, align contextual descriptions with time, or adapt predictions under event-driven changes~\cite{qiu2024tfb, aksu2024gifteval}. 
Conversely, temporal reasoning benchmarks for language models typically operate on purely textual data, omitting real numerical time-series signals~\cite{ning2020torque, fatemi2024test, wang2024tram}. 
As a result, it remains unclear whether agents genuinely improve forecasting behavior under context, or whether apparent gains stem from evaluation protocols that conflate numerical accuracy with broader capabilities.

To address this gap, we introduce \textbf{TemporalBench}, a multi-domain benchmark for evaluating agent and LLM performance on time-series tasks under varying levels of contextual and event information. 
TemporalBench spans four real-world domains, including retail, healthcare, energy, and physical systems, and is organized around a four-tier task taxonomy that jointly covers historical structure understanding, pure temporal prediction, contextual reasoning, and event-informed conditional forecasting. 
By selectively controlling access to future targets and contextual information, the benchmark disentangles complementary temporal competencies that cannot be assessed through numerical accuracy alone.

Beyond aggregate accuracy, TemporalBench serves as a diagnostic tool for revealing systematic failure modes of current agent systems, particularly in contextual grounding, temporal alignment, and robust use of historical information.

Overall, our contributions are as follows:
\begin{itemize}
\item We introduce \textbf{TemporalBench}, a multi-domain benchmark for evaluating agent and LLM performance on time-series tasks under contextual and event-informed settings.
\item We propose a four-tier task taxonomy (T1--T4) that disentangles historical understanding, pure temporal prediction, contextual reasoning, and event-conditioned forecasting.
\item We develop a unified benchmark construction pipeline with automatic task generation, rule-based ground-truth labeling, and explicit uncertainty modeling.
\item Through extensive baseline evaluation, we demonstrate that strong forecasting accuracy does not reliably translate into robust contextual or event-aware temporal reasoning.
\end{itemize}
